\documentclass[11pt,a4paper]{article}
% Shared preamble for the qi26_25 weekly deliverable reports.
%
% Each weekly report is a short standalone document that inputs this file, so
% that formatting stays consistent across all six without duplicating the
% package list in every one. Change something here and every report follows.
%
% Usage in a weekly report:
%     \documentclass[11pt,a4paper]{article}
%     % Shared preamble for the qi26_25 weekly deliverable reports.
%
% Each weekly report is a short standalone document that inputs this file, so
% that formatting stays consistent across all six without duplicating the
% package list in every one. Change something here and every report follows.
%
% Usage in a weekly report:
%     \documentclass[11pt,a4paper]{article}
%     % Shared preamble for the qi26_25 weekly deliverable reports.
%
% Each weekly report is a short standalone document that inputs this file, so
% that formatting stays consistent across all six without duplicating the
% package list in every one. Change something here and every report follows.
%
% Usage in a weekly report:
%     \documentclass[11pt,a4paper]{article}
%     \input{weekly_preamble}
%     \weekhead{1}{Vision Foundation and Data Encoding}{Week of ...}
%     ... body ...
%     \end{document}

\usepackage[utf8]{inputenc}
\usepackage[T1]{fontenc}
\usepackage[margin=2.5cm]{geometry}
\usepackage{amsmath,amssymb}
\usepackage{booktabs}
\usepackage{graphicx}
\usepackage{enumitem}
\usepackage{listings}
\usepackage{xcolor}
\usepackage[hidelinks]{hyperref}
\usepackage{titlesec}
\usepackage{fancyhdr}

\definecolor{codebg}{RGB}{247,247,250}
\definecolor{codekw}{RGB}{20,90,160}
\definecolor{codecm}{RGB}{110,120,130}
\definecolor{accent}{RGB}{20,90,160}

\lstset{
  backgroundcolor=\color{codebg},
  basicstyle=\ttfamily\footnotesize,
  keywordstyle=\color{codekw}\bfseries,
  commentstyle=\color{codecm}\itshape,
  breaklines=true,
  frame=single,
  rulecolor=\color{gray!30},
  language=Python,
  showstringspaces=false,
  xleftmargin=6pt,
  framexleftmargin=6pt
}

\setlist{itemsep=2pt,topsep=4pt}
\titleformat{\section}{\normalfont\large\bfseries\color{accent}}{\thesection}{0.8em}{}
\titleformat{\subsection}{\normalfont\normalsize\bfseries}{\thesubsection}{0.8em}{}

\pagestyle{fancy}
\fancyhf{}
\fancyfoot[C]{\small\thepage}
\fancyfoot[R]{\scriptsize\texttt{github.com/Pushkar0997/qi26\_25}}
\renewcommand{\headrulewidth}{0pt}

% \weekhead{number}{title}{deliverable as stated in the project brief}
\newcommand{\weekhead}[3]{%
  \begin{center}
    {\scriptsize\textsf{QINTERN 2026 \textbullet\ QWORLD \textbullet\ WEEKLY DELIVERABLE}}\\[0.25cm]
    {\LARGE\bfseries Week #1: #2}\\[0.35cm]
    {\normalsize Hybrid QML and Quantum String Algorithms for Document Extraction}\\[0.25cm]
    {\small Pushkar Kumar \quad\textbullet\quad Mentor: Potluri Krishna Priyatham}\\[0.4cm]
    \fbox{\begin{minipage}{0.88\textwidth}
      \small\textbf{Deliverable as specified in the project brief:} #3
    \end{minipage}}
  \end{center}
  \vspace{0.4cm}
}

% A boxed note for caveats and honest limitations.
\newcommand{\honestnote}[1]{%
  \vspace{0.2cm}
  \noindent\fcolorbox{gray!40}{gray!8}{%
    \begin{minipage}{0.97\textwidth}\small #1\end{minipage}}
  \vspace{0.2cm}
}

%     \weekhead{1}{Vision Foundation and Data Encoding}{Week of ...}
%     ... body ...
%     \end{document}

\usepackage[utf8]{inputenc}
\usepackage[T1]{fontenc}
\usepackage[margin=2.5cm]{geometry}
\usepackage{amsmath,amssymb}
\usepackage{booktabs}
\usepackage{graphicx}
\usepackage{enumitem}
\usepackage{listings}
\usepackage{xcolor}
\usepackage[hidelinks]{hyperref}
\usepackage{titlesec}
\usepackage{fancyhdr}

\definecolor{codebg}{RGB}{247,247,250}
\definecolor{codekw}{RGB}{20,90,160}
\definecolor{codecm}{RGB}{110,120,130}
\definecolor{accent}{RGB}{20,90,160}

\lstset{
  backgroundcolor=\color{codebg},
  basicstyle=\ttfamily\footnotesize,
  keywordstyle=\color{codekw}\bfseries,
  commentstyle=\color{codecm}\itshape,
  breaklines=true,
  frame=single,
  rulecolor=\color{gray!30},
  language=Python,
  showstringspaces=false,
  xleftmargin=6pt,
  framexleftmargin=6pt
}

\setlist{itemsep=2pt,topsep=4pt}
\titleformat{\section}{\normalfont\large\bfseries\color{accent}}{\thesection}{0.8em}{}
\titleformat{\subsection}{\normalfont\normalsize\bfseries}{\thesubsection}{0.8em}{}

\pagestyle{fancy}
\fancyhf{}
\fancyfoot[C]{\small\thepage}
\fancyfoot[R]{\scriptsize\texttt{github.com/Pushkar0997/qi26\_25}}
\renewcommand{\headrulewidth}{0pt}

% \weekhead{number}{title}{deliverable as stated in the project brief}
\newcommand{\weekhead}[3]{%
  \begin{center}
    {\scriptsize\textsf{QINTERN 2026 \textbullet\ QWORLD \textbullet\ WEEKLY DELIVERABLE}}\\[0.25cm]
    {\LARGE\bfseries Week #1: #2}\\[0.35cm]
    {\normalsize Hybrid QML and Quantum String Algorithms for Document Extraction}\\[0.25cm]
    {\small Pushkar Kumar \quad\textbullet\quad Mentor: Potluri Krishna Priyatham}\\[0.4cm]
    \fbox{\begin{minipage}{0.88\textwidth}
      \small\textbf{Deliverable as specified in the project brief:} #3
    \end{minipage}}
  \end{center}
  \vspace{0.4cm}
}

% A boxed note for caveats and honest limitations.
\newcommand{\honestnote}[1]{%
  \vspace{0.2cm}
  \noindent\fcolorbox{gray!40}{gray!8}{%
    \begin{minipage}{0.97\textwidth}\small #1\end{minipage}}
  \vspace{0.2cm}
}

%     \weekhead{1}{Vision Foundation and Data Encoding}{Week of ...}
%     ... body ...
%     \end{document}

\usepackage[utf8]{inputenc}
\usepackage[T1]{fontenc}
\usepackage[margin=2.5cm]{geometry}
\usepackage{amsmath,amssymb}
\usepackage{booktabs}
\usepackage{graphicx}
\usepackage{enumitem}
\usepackage{listings}
\usepackage{xcolor}
\usepackage[hidelinks]{hyperref}
\usepackage{titlesec}
\usepackage{fancyhdr}

\definecolor{codebg}{RGB}{247,247,250}
\definecolor{codekw}{RGB}{20,90,160}
\definecolor{codecm}{RGB}{110,120,130}
\definecolor{accent}{RGB}{20,90,160}

\lstset{
  backgroundcolor=\color{codebg},
  basicstyle=\ttfamily\footnotesize,
  keywordstyle=\color{codekw}\bfseries,
  commentstyle=\color{codecm}\itshape,
  breaklines=true,
  frame=single,
  rulecolor=\color{gray!30},
  language=Python,
  showstringspaces=false,
  xleftmargin=6pt,
  framexleftmargin=6pt
}

\setlist{itemsep=2pt,topsep=4pt}
\titleformat{\section}{\normalfont\large\bfseries\color{accent}}{\thesection}{0.8em}{}
\titleformat{\subsection}{\normalfont\normalsize\bfseries}{\thesubsection}{0.8em}{}

\pagestyle{fancy}
\fancyhf{}
\fancyfoot[C]{\small\thepage}
\fancyfoot[R]{\scriptsize\texttt{github.com/Pushkar0997/qi26\_25}}
\renewcommand{\headrulewidth}{0pt}

% \weekhead{number}{title}{deliverable as stated in the project brief}
\newcommand{\weekhead}[3]{%
  \begin{center}
    {\scriptsize\textsf{QINTERN 2026 \textbullet\ QWORLD \textbullet\ WEEKLY DELIVERABLE}}\\[0.25cm]
    {\LARGE\bfseries Week #1: #2}\\[0.35cm]
    {\normalsize Hybrid QML and Quantum String Algorithms for Document Extraction}\\[0.25cm]
    {\small Pushkar Kumar \quad\textbullet\quad Mentor: Potluri Krishna Priyatham}\\[0.4cm]
    \fbox{\begin{minipage}{0.88\textwidth}
      \small\textbf{Deliverable as specified in the project brief:} #3
    \end{minipage}}
  \end{center}
  \vspace{0.4cm}
}

% A boxed note for caveats and honest limitations.
\newcommand{\honestnote}[1]{%
  \vspace{0.2cm}
  \noindent\fcolorbox{gray!40}{gray!8}{%
    \begin{minipage}{0.97\textwidth}\small #1\end{minipage}}
  \vspace{0.2cm}
}


\begin{document}
\weekhead{1}{Vision Foundation and Data Encoding}
{Technical summary of encoding efficiency versus classical normalisation.}

\section{Objective}

Establish a quantum representation for document imagery. Two standard encodings
were implemented and compared: FRQI (Flexible Representation of Quantum Images)
and NEQR (Novel Enhanced Quantum Representation). The question to be answered is
which is appropriate for a document-processing pipeline, and at what resource
cost relative to simply normalising pixels classically.

\section{The Two Encodings}

\textbf{FRQI} stores each pixel's intensity as a rotation angle on a single
shared colour qubit:
\begin{equation}
|I\rangle = \frac{1}{2^{n}}\sum_{i=0}^{2^{2n}-1}
\big(\cos\theta_i\,|0\rangle + \sin\theta_i\,|1\rangle\big)\otimes|i\rangle,
\qquad \theta_i = \frac{\pi}{2}\cdot\frac{v_i}{255}
\end{equation}

One qubit carries all colour information regardless of image size. The intensity
range maps to $[0, \pi/2]$ rather than $[0,\pi]$ so that $\sin^2$ remains
monotonic and a measured probability inverts to exactly one intensity.

\textbf{NEQR} stores intensity as a bit string on a colour register:
\begin{equation}
|I\rangle = \frac{1}{2^{n}}\sum_{i}|f(i)\rangle\otimes|i\rangle
\end{equation}

An 8-bit greyscale image requires 8 colour qubits instead of 1, but
reconstruction is exact.

\section{Verification}

A circuit that runs is not a circuit that is correct, so both encodings were
inverted and checked against the original pixel values rather than assumed
correct.

For FRQI, the colour qubit was measured conditioned on each position and the
encoding inverted:
\begin{equation}
v_i = \frac{255 \cdot 2}{\pi}\arcsin\sqrt{P(\text{colour}=1 \mid i)}
\end{equation}

\begin{table}[h]
\centering
\caption{Reconstruction accuracy on a $2\times2$ patch.}
\begin{tabular}{@{}lll@{}}
\toprule
Encoding & Reconstruction & Maximum error \\
\midrule
FRQI & statistical, from measurement frequencies & 1--2 intensity levels of 255 \\
NEQR & direct bit readout & \textbf{exactly 0} \\
\bottomrule
\end{tabular}
\end{table}

The FRQI residual is sampling error and shrinks as $1/\sqrt{\text{shots}}$; the
figure above is at 8{,}192 shots. NEQR's error is exactly zero because no
estimation is involved.

\section{Resource Comparison}

Both circuits were transpiled to a common $\{u, \mathrm{CX}\}$ basis before
counting. This matters for fairness: multi-controlled gates are single
instructions at the abstract level but expand into very different numbers of
two-qubit gates depending on their control count, so counting before
transpilation would make a six-control gate appear as cheap as a CX.

\begin{table}[h]
\centering
\caption{Measured resources across three patch sizes.}
\begin{tabular}{@{}lrrrrrr@{}}
\toprule
Patch & Pixels & FRQI qubits & NEQR qubits & FRQI CX & NEQR CX & Qubit ratio \\
\midrule
$2\times2$ & 4 & 3 & 10 & 32 & 114 & $3.33\times$ \\
$4\times4$ & 16 & 5 & 12 & 384 & 1{,}558 & $2.40\times$ \\
$8\times8$ & 64 & 7 & 14 & 3{,}584 & 10{,}998 & $2.00\times$ \\
\bottomrule
\end{tabular}
\end{table}

\section{A Prediction, and Its Refutation}

The initial version of this summary, written before the $4\times4$ and $8\times8$
measurements were taken, predicted that the qubit gap between the two encodings
would \textbf{widen} as patch size grew.

Measurement refutes this. The ratio \textbf{narrows}: $3.33\times \rightarrow
2.40\times \rightarrow 2.00\times$.

The original reasoning conflated two distinct quantities. FRQI's colour register
is fixed at 1 qubit and NEQR's at 8, so the colour gap is a \emph{constant} 7
qubits at every patch size. Meanwhile both encodings share an identically growing
position register of $\lceil\log_2(\text{pixels})\rceil$. A constant gap divided
by a growing shared base necessarily shrinks as a ratio.

\honestnote{%
The prediction was wrong and is recorded here rather than quietly corrected.
Absolute CX counts do grow steeply for both encodings---NEQR reaches
approximately 11{,}000 gates at $8\times8$---so NEQR remains far more expensive
in absolute terms. It is specifically the \emph{relative} qubit penalty that
decays.}

\section{Comparison Against Classical Normalisation}

The brief asked for encoding efficiency relative to classical normalisation. The
comparison is stark and worth stating plainly.

Classical normalisation of a $2\times2$ patch is four floating-point divisions:
approximately four operations, exact, with no measurement required. FRQI requires
3 qubits and 32 CX gates and returns an \emph{approximate} value needing
thousands of shots to estimate. NEQR requires 10 qubits and 114 CX gates for an
exact value.

No efficiency argument favours quantum encoding at this scale. The encodings are
of interest only as a route into subsequent quantum processing---if the
subsequent processing provides an advantage, which is the question Weeks 2 and 4
address.

\section{Consequence for the Pipeline}

The pipeline ultimately uses \emph{neither} encoding directly, and it is worth
being explicit about why rather than leaving the reader to infer it.

The quanvolutional layer developed in Week 2 uses a simple $R_Y$ angle encoding:
one qubit per pixel of a $2\times2$ patch, rotated by $\pi v/255$. This is
FRQI-like in spirit---intensity as angle---but has no position register, because
a quanvolutional filter processes one patch at a time and never needs to address
pixels by index.

FRQI and NEQR remain relevant as the resource baseline for \emph{whole-image}
quantum representation, which any approach operating on a full document at once
would require. The measurement that NEQR costs roughly 11{,}000 CX gates for a
single $8\times8$ patch is the quantitative reason this pipeline is patch-based
rather than whole-image.

\section{Deliverable Status}

Complete. Both encodings implemented and verified by inversion; resource
comparison measured at three patch sizes; the hypothesis about scaling tested and
refuted.

\vspace{0.3cm}
\noindent\textbf{Artifacts:} \texttt{track\_a\_vision/encoding/encoding\_scaling.py},
\texttt{notebooks/01\_encoding\_frqi\_neqr.ipynb},
\texttt{benchmarking/encoding\_scaling.json}

\end{document}
